\section{What to pursue next}
There is no single well-defined ``simplest'' polynomial without a cost
model. The present evidence supports the following priorities rather
than a claim that the target has been minimized.

\paragraph{Compile the height source before optimizing its presentation.}
The largest architectural opportunity is the removal of an entire
bounded-universal common-multiple certificate. Instantiate and elaborate
the supplied constructor, prove the local relation's semantics in the
chosen kernel, and measure the full bounded-forall and remainder expansion.
Only then can the extra accumulator stream be weighed against the deleted
block. The 23/14 local count cannot predict the global variable count by
itself. This is the most direct test of whether the structural idea beats
the Pell-based endpoints.

\paragraph{Audit how much Lucas divisibility is really needed.}
Norm saturation succeeds because $D$ becomes a square in the soundness
proof, while $\Omega$ still absorbs the relevant period valuations.
Removing another factor from $\Omega$ or $C^2$ would require a new rank
and full-residue interpolation theorem. The counterexample at $C=-1$
shows why a locally attractive formula cannot be accepted by gate counting
alone. The correct research object is the exact divisibility interface,
including exceptional indices, not an unchecked symbolic cancellation.

\paragraph{Reduce the source while keeping complete coefficient tuples.}
Compatible packing can lower degree without changing witnesses, and the
factored positive gap can remove a weighted source record. These are
separate levers. Any proposed replacement must preserve all extraction
identities, carry bounds, positivity, canonical divisibilities and the
whole fixed tuple. A lower syntactic term count without a better
certified layout is not an improved final equation.

\paragraph{Keep degree and cost tradeoffs honest.}
The new sextic uses retained $F_2=f^2$ as an ideal-preserving lifting
register. Other registers may admit similar cancellations, but any
elimination must retain integrality and a polynomial restoration map.
For the positive equations, rational resultants need a descent or
saturation argument; elimination over $\Q$ alone does not establish a
Diophantine equivalence over $\Z$.

